% =============================================================================
% 6 Plausibilisierung und Bewertung
% =============================================================================
\section{Plausibilisierung und Bewertung}
\label{chap:plausibilisierung_bewertung}

% -----------------------------------------------------------------------------
% 6.1 Vorgehen und untersuchte Vergleichsfälle
% -----------------------------------------------------------------------------
\subsection{Vorgehen und untersuchte Vergleichsfälle}
\label{sec:vorgehen_vergleichsfaelle}

% Inhalt:
% - Ziel der Plausibilisierung
% - Abgrenzung gegenüber einer vollständigen Validierung
% - erwartetes physikalisches Verhalten
% - Festlegung der untersuchten Grenz- und Vergleichsfälle
%
% Vergleichsfälle:
% - Starttemperatur 20 °C
% - Starttemperatur 30 °C
% - P_el = 0
% - konstante Verlustleistung
% - Spannungssprung
% - stationärer Endzustand
% - analytische PT1-Lösung


% -----------------------------------------------------------------------------
% 6.2 Plausibilisierung der Teilmodelle
% -----------------------------------------------------------------------------
\subsection{Plausibilisierung der Teilmodelle}
\label{sec:plausibilisierung_teilmodelle}

% Inhalt des thermischen Teilmodells:
% - Abkühlung beziehungsweise unveränderte Temperatur bei P_el = 0
% - Erwärmung bei konstanter Verlustleistung
% - Einfluss unterschiedlicher Anfangstemperaturen
% - stationäre Endtemperatur
% - Vergleich mit der analytischen PT1-Lösung
%
% Inhalt des elektrischen Teilmodells:
% - dreiphasiger Spannungsverlauf
% - Effektivwertbildung
% - Reaktion auf den Spannungssprung
% - Strom bei konstantem Widerstand
% - Einfluss von R_Stell
% - Stromaufteilung im Zusatzpfad


% -----------------------------------------------------------------------------
% 6.3 Verhalten ohne Stromnachführung
% -----------------------------------------------------------------------------
\subsection{Verhalten ohne Stromnachführung}
\label{sec:verhalten_ohne_stromnachfuehrung}

% Inhalt:
% - Temperaturverlauf des V2A-Widerstands
% - Widerstandsanstieg
% - Abfall des Modulstroms
% - Veränderung der Verlustleistung
% - stationäres Verhalten
% - Nachbildung der bestehenden Problematik


% -----------------------------------------------------------------------------
% 6.4 Verhalten mit Stromnachführung
% -----------------------------------------------------------------------------
\subsection{Verhalten mit Stromnachführung}
\label{sec:verhalten_mit_stromnachfuehrung}

% Inhalt:
% - Aktivierung des Stellpfads
% - Veränderung von R_Stell
% - Nachführung des Modulstroms
% - erreichbare Regelgenauigkeit
% - erforderlicher Stellweg
% - verfügbare Stellreserve
% - gegenseitige Beeinflussung der Abgangseinheiten


% -----------------------------------------------------------------------------
% 6.5 Bewertung des Optimierungskonzepts
% -----------------------------------------------------------------------------
\subsection{Bewertung des Optimierungskonzepts}
\label{sec:bewertung_optimierungskonzept}

% Inhalt:
% - Vergleich mit den Anforderungen aus Kapitel 3
% - individuelle Stromnachführung
% - Mehrkanalfähigkeit
% - Stellbereich und Stellreserve
% - Dauerbetriebsfähigkeit
% - Integration in den bestehenden Prüfstand
% - erreichbare Verbesserung gegenüber dem Ist-Zustand
%
% Darstellung möglichst als Tabelle:
% Anforderung | Ergebnis | erfüllt / teilweise erfüllt / nicht nachgewiesen


% -----------------------------------------------------------------------------
% 6.6 Modellgrenzen und Unsicherheiten
% -----------------------------------------------------------------------------
\subsection{Modellgrenzen und Unsicherheiten}
\label{sec:modellgrenzen_unsicherheiten}

% Inhalt:
% - Annahme eines konstanten Wärmeübergangs
% - vereinfachte thermische Parameter
% - reduzierte Abbildung der Transformatoren
% - vernachlässigte Kontakt- und Leitungswiderstände
% - nicht vollständig abgebildete Kopplungen
% - fehlender beziehungsweise begrenzter Messwertvergleich
% - Grenzen der Übertragbarkeit auf den realen Prüfstand